\documentclass[10pt,a4paper]{amsart}
\usepackage[margin=19mm]{geometry}
\usepackage{amsmath,amssymb,mathtools}
\usepackage[scr=rsfs,cal=euler]{mathalfa}
\usepackage{xfrac}
\usepackage[unicode,hidelinks,bookmarks=false]{hyperref}
\newcommand{\ie}{i.e.\ }
\newcommand{\cf}{cf.\ }
\newcommand{\resp}{respectively\ }
\newcommand{\Subsection}{\subsection}
\providecommand{\nobreakdash}{\nobreak\hbox{-}}
\setlength{\emergencystretch}{2em}
\multlinegap0pt
\usepackage{amssymb}
\usepackage{cancel}
\newtheorem{lemma}{Lemma}
\newcommand{\CH}{\mathsf{CH}}
\newcommand{\forceA}{\mathbb{A}}
\newcommand{\forceC}{\mathbb{C}}
\newcommand{\forceP}{\mathbb{P}}
\newcommand{\forceQ}{\mathbb{Q}}
\newcommand{\forceR}{\mathbb{R}}
\newcommand{\forceS}{\mathbb{S}}
\newcommand{\key}{\operatorname{key}}
\title{Forcing the ${\Pi^1_n}$-Uniformization Property}
\author{Stefan Hoffelner}
\date{Source: arXiv:2103.11748v5 (CC BY 4.0)}
\begin{document}
\maketitle

\noindent\emph{Source and licence.} Forcing the ${\Pi^1_n}$-Uniformization Property. Version arXiv:2103.11748v5 (CC BY 4.0), distributed by arXiv under the Creative Commons Attribution 4.0 International licence, http://creativecommons.org/licenses/by/4.0/, as recorded in the arXiv licence metadata for this exact version and shown on the abstract page https://arxiv.org/abs/2103.11748v5. Full licence notice: \texttt{licenses/arxiv-2103.11748-CC-BY-4.0.txt}.\par\smallskip

\noindent\textit{Editorial scope.} Counted unit: the article's lemma 1allowableclosedunderproducts (source0.tex lines 1830--1844). The article's complete original proof of it is delivered with it (source0.tex lines 1846--1973, one \texttt{proof} environment of 128 lines). No mathematics is written, repaired or re-derived: the statement and the proof are the article's own lines, in the article's own notation, and no retained line is substituted or re-flowed.  Retained uncounted as the source-local context the proof needs: lines 805--832; lines 1773--1781.  Nothing was omitted from any retained span, so the package carries no omission record.  No retained line contains a citation, so the wrapper carries no bibliography and no bibliography entry.  No retained line contains a figure environment, an image-inclusion command or drawing code, so no image is delivered and the package declares no asset dependency.  Editorial formatting only, in addition: an AMS \texttt{amsart} preamble that restates the article's own declarations and macros used by the retained lines, namely the environments \texttt{lemma} on separate counters, together with the standard packages those lines need.  The retained environments print in this document as Lemma 1 in order of appearance, whereas the article prints the same items with the numbering of its own full document.  The complete native source of the article as distributed in the arXiv e-print for this exact version is bundled unchanged as \texttt{sources/109407/source0.tex}.
\begin{lemma}
\label{basic-properties-allowable-forcings}
Let $\forceP_\delta=\forceC^{\mathrm{area}}*\dot{\forceR}_\delta$ be an
allowable forcing of length $\delta<\omega_1$.
\begin{enumerate}
   \item For every $\beta<\delta$,
   \[
      \forceP_\beta\Vdash
      ``\dot{\forceA}_\beta\text{ is Knaster and }
      |\dot{\forceA}_\beta|=\aleph_1."
   \]
   Consequently,
   $\forceC^{\mathrm{area}}\Vdash
   ``\dot{\forceR}_\delta\text{ is c.c.c. and has size at most }
   \aleph_1."$

   \item $\forceP_\delta$ preserves $\omega_1$ and $\CH$.

   \item If $\forceP^0$ and $\forceP^1$ are allowable, then
   \[
      \forceP^0\times\forceP^1
   \]
   is isomorphic to an allowable forcing.  More precisely, after an
   $L$-definable reindexing of the two reservoir factors onto disjoint
   sets of coordinates, the product admits an allowable iteration over the
   global reservoir.
\end{enumerate}
\end{lemma}

\begin{lemma}
\label{shrinkinglemma}
Let $\xi<\alpha$.  Every $\alpha$-allowable forcing is $\xi$-allowable.  Consequently,
\[
   \mathcal A^W_0\supseteq\mathcal A^W_1\supseteq\cdots
   \supseteq\mathcal A^W_\alpha\supseteq\cdots,
\]
where $\mathcal A^W_\alpha$ denotes the class of $\alpha$-allowable forcings over $W$.
\end{lemma}

\begin{lemma}
\label{1allowableclosedunderproducts}
\label{alphaallowableclosedunderproducts}
Let $\alpha$ be an ordinal.  If $\forceP^0$ and $\forceP^1$ are
$\alpha$-allowable, then
\[
   \forceP^0\times\forceP^1
\]
is isomorphic to an $\alpha$-allowable forcing.  More precisely, a normalized
$\alpha$-allowable presentation of the product is obtained from the two given
presentations by an $L$-definable reindexing of the two reservoirs onto
disjoint blocks and by concatenating the translated almost disjoint tails.  The
resulting bookkeeping function belongs to $W$, and the resulting injective area
assignment belongs to $L$.
\end{lemma}

\begin{proof}
We argue by induction on $\alpha$, proving the normalized form stated above.
For $\alpha=0$, this is exactly
Lemma~\ref{basic-properties-allowable-forcings}(3).

Assume that the assertion holds for every $\xi<\alpha$.  Let
$(F_i,\vec\eta^{\,i})$ witness the $\alpha$-allowability of
$\forceP^i$, and let $\delta_i$ be the length of $\forceP^i$
$(i<2)$.  Fix in $L$ a partition
\[
   \omega_1=B_0\mathbin{\dot\cup}B_1
\]
and bijections
\[
   e_i:\omega_1\longrightarrow B_i
   \qquad(i<2).
\]
Move the reservoir of $\forceP^i$ to $B_i$ by $e_i$.  We present the product
by first taking the translated almost disjoint tail of $\forceP^0$ and then
the translated tail of $\forceP^1$.  Let
\[
   \forceS=\langle\forceS_\gamma\mid
      \gamma\leq\delta_0+\delta_1\rangle
\]
be this normalized product presentation.  Its final forcing is isomorphic to
$\forceP^0\times\forceP^1$.  Put
\[
   \vec\eta^{\,\otimes}
   =
   \langle e_0(\eta^{\,0}_\beta)\mid\beta<\delta_0\rangle
    {}^\frown
    \langle e_1(\eta^{\,1}_\beta)\mid\beta<\delta_1\rangle.
\]
Then $\vec\eta^{\,\otimes}\in L$ is injective.  Let
$F^\otimes$ be obtained by translating $F_0$ and $F_1$ by $e_0$ and
$e_1$, respectively, and concatenating the translated sequences.  Then
$F^\otimes\in W$.

It remains to check the stage rule.  The first block is handled by the same
argument as below, with only the reindexing $e_0$, so we consider a stage
$\delta_0+\beta$ in the second block.  Let $\beta<\delta_1$, and let
\[
   \mathfrak a^1_\beta=(E,A)
\]
be the local admissible support of the corresponding stage in the construction
of $\forceP^1$.  In the product presentation the translated support is
\[
   \mathfrak a^\otimes_{\delta_0+\beta}
   :=(e_1[E],\delta_0+A),
   \qquad
   \delta_0+A=\{\delta_0+\gamma:\gamma\in A\}.
\]
It is admissible, since
\[
   E\subseteq
   \operatorname{ran}(\vec\eta^{\,1}\upharpoonright\beta)
\]
implies
\[
   e_1[E]\subseteq
   \operatorname{ran}
      (\vec\eta^{\,\otimes}\upharpoonright(\delta_0+\beta)).
\]
The local models
\[
   W[K\upharpoonright e_1[E]]
    [G\upharpoonright(\delta_0+A)]
   \quad\text{and}\quad
   W[K^1\upharpoonright E][G^1\upharpoonright A]
\]
are canonically isomorphic.  We identify their reals through this isomorphism.

Fix $\xi<\alpha$ and a local real $y$.  We claim that $y$ is
$\xi$-persistent for $(x,m)$ at the stage $\beta$ of $\forceP^1$ iff its
translated copy is $\xi$-persistent at the stage $\delta_0+\beta$ of
$\forceS$.  If $y$ is not persistent in $\forceP^1$, let $\forceQ$ be a
$\xi$-allowable complete reservoir extension of $\forceP^1_\beta$ whose
quotient forces
\[
   (x,y)\notin A_m.
\]
By Lemma~\ref{shrinkinglemma}, $\forceP^0$ is $\xi$-allowable, and by the
induction hypothesis $\forceP^0\times\forceQ$ has a normalized
$\xi$-allowable presentation.  This presentation contains
$\forceS_{\delta_0+\beta}$ as a complete reservoir subiteration, and its
quotient still forces
\[
   (x,y)\notin A_m.
\]
Hence the translated $y$ is not persistent in the product presentation.

Conversely, if a $\xi$-allowable complete reservoir extension of
$\forceS_{\delta_0+\beta}$ forces
\[
   (x,y)\notin A_m,
\]
then, since $\forceS_{\delta_0+\beta}$ contains the translated copy of
$\forceP^1_\beta$ as a complete reservoir subiteration, the same forcing is
also a $\xi$-allowable complete reservoir extension of $\forceP^1_\beta$ by
transitivity of $\succeq_{\mathrm{res}}$.  Thus $y$ is not persistent in the
original construction of $\forceP^1$.  This proves the claim.

The key comparison is also preserved.  The reindexing $e_1$ sends presentations
below $\mathfrak a^1_\beta$ bijectively to presentations below
$\mathfrak a^\otimes_{\delta_0+\beta}$, and the corresponding pointed
presentations have the same presentation class.  Hence, for every local real
$y$,
\[
   \key_{\mathfrak a^1_\beta}
      (\langle x,y,m\rangle)
   =
   \key_{\mathfrak a^\otimes_{\delta_0+\beta}}
      (\langle x,y,m\rangle).
\]
Equivalently, the two constructions compute the same orbits and therefore the
same $<_L$-least keys.  Thus they have the same least persistence rank and,
at that rank, choose the same real.  If the selected value equals the
bookkeeping value, the same key comparison also chooses the same real to be
coded instead.  Therefore the translated stage of $\forceS$ uses exactly the
translated coding factor from $\forceP^1$.

The same verification applies to every stage of the second block, and the
first block was already covered by the reindexing argument.  At limits the
finite-support limits agree.  Hence $\forceS$ is an $\alpha$-allowable
iteration relative to $(F^\otimes,\vec\eta^{\,\otimes})$.  Since the final
forcing of $\forceS$ is isomorphic to $\forceP^0\times\forceP^1$, the product
is isomorphic to an $\alpha$-allowable forcing.
\end{proof}

\end{document}
